
\subsection*{\SGAInlineRefTarget{sga3:i:subsection:2-3}2.3. Structures with operator groups}

\medskip
\noindent\SGAInlineRefTarget{sga3:i:definition:2-3-1}\textbf{Definition 2.3.1.}
Let \(\mathbf E\in\operatorname{Ob}(\widehat{\mathscr C})\) and
\(\mathbf G\in
\operatorname{Ob}(\widehat{\mathscr C}\text{-}\mathbf{Gr})\).
A structure on \(\mathbf E\) of an object with
\(\widehat{\mathscr C}\)-operator group \(\mathbf G\) (or of a
\(\mathbf G\)-object) consists in giving \(\mathbf E(S)\), for every
\(S\in\operatorname{Ob}(\mathscr C)\), the structure of a set with
operator group \(\mathbf G(S)\), in such a way that, for every arrow
\(S'\to S''\) in \(\mathscr C\), the map of sets
\(\mathbf E(S'')\to\mathbf E(S')\) is compatible with the homomorphism
of operator groups \(\mathbf G(S'')\to\mathbf G(S')\).

As usual, this is equivalent to giving a morphism
\[
  \mu:\mathbf G\times\mathbf E\longrightarrow\mathbf E
\]
which, for every \(S\), equips \(\mathbf E(S)\) with the structure of a
set on which \(\mathbf G(S)\) acts. But
\[
  \operatorname{Hom}(\mathbf G\times\mathbf E,\mathbf E)
    \simeq
  \operatorname{Hom}
    (\mathbf G,\underline{\operatorname{End}}(\mathbf E)),
\]
so \(\mu\) defines a morphism
\(\mathbf G\to\underline{\operatorname{End}}(\mathbf E)\). It is
immediate that this morphism maps \(\mathbf G\) into
\(\underline{\operatorname{Aut}}(\mathbf E)\) and is a morphism of
\(\widehat{\mathscr C}\)-groups. Consequently, giving \(\mathbf E\) the
structure of an object with \(\widehat{\mathscr C}\)-operator group
\(\mathbf G\) is equivalent to giving a morphism of
\(\widehat{\mathscr C}\)-groups
\[
  \rho:\mathbf G\longrightarrow
    \underline{\operatorname{Aut}}(\mathbf E).
\]
In particular, every element \(g\in\mathbf G(S)\) defines an
automorphism \(\rho(g)\) of the functor \(\mathbf E_S\), that is, an
automorphism of \(\mathbf E\times\mathbf h_S\) commuting with the
projection \(\mathbf E\times\mathbf h_S\to\mathbf h_S\), and hence, in
particular, an automorphism of the set \(\mathbf E(S')\) for every
\(S'\to S\).

\medskip
\noindent\SGAInlineRefTarget{sga3:i:definition:2-3-2}\textbf{Definition 2.3.2.}
We denote by \(\mathbf E^{\mathbf G}\) the subobject of \(\mathbf E\)
defined as follows:
\[
  \mathbf E^{\mathbf G}(S)
    =\bigl\{x\in\mathbf E(S)\mid
      x_{S'}\text{ is invariant under }\mathbf G(S')
      \text{ for every }S'\to S\bigr\},
\]
where \(x_{S'}\) denotes the image of \(x\) under
\(\mathbf E(S)\to\mathbf E(S')\).

Thus \(\mathbf E^{\mathbf G}\), the ``subobject of invariants of
\(\mathbf G\),'' is the largest subobject of \(\mathbf E\) on which
\(\mathbf G\) acts trivially.

\medskip
\SGARefTarget{sga3:i:definition:2:3:3}
\noindent\textbf{Definition 2.3.3.}
Let \(\mathbf F\) be a subobject of \(\mathbf E\). We denote by
\(\underline{\operatorname{Norm}}_{\mathbf G}\mathbf F\) and
\(\underline{\operatorname{Centr}}_{\mathbf G}\mathbf F\) the
sub-\(\widehat{\mathscr C}\)-groups of \(\mathbf G\) defined by
\begin{align*}
  (\underline{\operatorname{Norm}}_{\mathbf G}\mathbf F)(S)
    &={}
      \{g\in\mathbf G(S)\mid \rho(g)\mathbf F_S=\mathbf F_S\}\\
    &={}
      \{g\in\mathbf G(S)\mid
        \rho(g)\mathbf F(S')=\mathbf F(S')
        \text{ for every }S'\to S\},\\
  (\underline{\operatorname{Centr}}_{\mathbf G}\mathbf F)(S)
    &={}
      \{g\in\mathbf G(S)\mid
        \rho(g)|_{\mathbf F_S}=\operatorname{id}\}\\
    &={}
      \{g\in\mathbf G(S)\mid
        \rho(g)|_{\mathbf F(S')}=\operatorname{id}
        \text{ for every }S'\to S\}.
\end{align*}
Here the vertical bar after \(\rho(g)\) denotes restriction.\footnote{Translator's note: in the definition of
\(\underline{\operatorname{Norm}}_{\mathbf G}\mathbf F\), the 1970 text has
the inclusions \(\rho(g)\mathbf F_S\subset\mathbf F_S\) and
\(\rho(g)\mathbf F(S')\subset\mathbf F(S')\); the Gille--Polo re-edition
(2011), Exposé~I, replaces them by equalities, so that one obtains a
subgroup. With the inclusions, the set of such \(g\) is stable under
products but not always under inverses: let the constant group
\(S\mapsto\mathbb Z\) act by translations on the constant functor
\(S\mapsto\mathbb Z\), and let \(\mathbf F\) be the constant subfunctor
\(S\mapsto\mathbb N\); then \(1\) maps \(\mathbf F\) into itself and \(-1\)
does not. See Exposé~VI\textsubscript{B}, \SGARefLink{prop:sga3-VIB-6.4}{6.4}, for conditions under
which the two definitions agree.}

\par\medskip\noindent\SGAInlineRefTarget{sga3:i:scholium:2-3-3-1}%
\emph{Addition} (after the Gille--Polo re-edition (2011), Exposé~I, 2.3.3.1
and 2.3.3.2; our formulation). A section \(x\in\Gamma(\mathbf E)\) (1.2),
that is, a compatible family of elements \(x_S\in\mathbf E(S)\) (an element
of \(\mathbf E(S_0)\) when \(\mathscr C\) has a final object \(S_0\)),
defines the subobject \(S\mapsto\{x_S\}\) of \(\mathbf E\); its normalizer
and its centralizer coincide, and this sub-\(\widehat{\mathscr C}\)-group,
the \emph{stabilizer} \(\underline{\operatorname{Stab}}_{\mathbf G}(x)\) of
\(x\), is given by
\(\underline{\operatorname{Stab}}_{\mathbf G}(x)(S)=\{g\in\mathbf G(S)\mid\rho(g)x_S=x_S\}\).
If \(\mathscr C\) has a final object \(S_0\) and fibre products,
\(\mathbf G=\mathbf h_G\) and \(\mathbf E=\mathbf h_E\), then
\(x:S_0\to E\), and \(\underline{\operatorname{Stab}}_{\mathbf G}(x)\) is
represented by the fibre product \(G\times_ES_0\) of \(x\) and of the
morphism \(G\to E\), \(g\mapsto gx\), composite of
\(\operatorname{id}_G\times x:G\to G\times E\) and of the action. Moreover,
the formation of \(\mathbf E^{\mathbf G}\),
\(\underline{\operatorname{Norm}}_{\mathbf G}\mathbf F\) and
\(\underline{\operatorname{Centr}}_{\mathbf G}\mathbf F\) commutes with base
change: \((\mathbf E^{\mathbf G})_S=(\mathbf E_S)^{\mathbf G_S}\), and
similarly for the other two, their values at an object \(S'\to S\) only
involving the objects over \(S'\).
\par\medskip

\medskip
\noindent\SGAInlineRefTarget{sga3:i:definition:2-3-4}\textbf{Definition 2.3.4.}
If \(G\) is a \(\mathscr C\)-group and \(\mathbf E\) is an object of
\(\widehat{\mathscr C}\) (respectively, \(E\) is an object of
\(\mathscr C\)), a \(G\)-object structure on \(\mathbf E\)
(respectively, on \(E\)) is an \(\mathbf h_G\)-object structure on
\(\mathbf E\) (respectively, on \(\mathbf h_E\)).

In view of this definition, all the notions and notation defined above
carry over to \(\mathscr C\) whenever they involve only representable
functors. For example, if
\(\underline{\operatorname{Norm}}_{\mathbf h_G}(\mathbf h_F)\) is
representable, then there is a unique subobject of \(G\) representing
it; it is then a sub-\(\mathscr C\)-group of \(G\), and is denoted by
\(\operatorname{Norm}_G(F)\), and so forth.

\medskip
\noindent\SGAInlineRefTarget{sga3:i:definition:2-3-5}\textbf{Definition 2.3.5.}
\begin{enumerate}
  \item[(a)]
  We say that the \(\widehat{\mathscr C}\)-group \(\mathbf G\) acts on
  the \(\widehat{\mathscr C}\)-group \(\mathbf H\) if \(\mathbf H\)
  is equipped with a \(\mathbf G\)-object structure such that, for
  every \(g\in\mathbf G(S)\), the automorphism of \(\mathbf H(S)\)
  defined by \(g\) is a group automorphism.

  Equivalently, for every \(g\in\mathbf G(S)\), the automorphism
  \(\rho(g)\) of \(\mathbf H_S\) is an automorphism of
  \(\widehat{\mathscr C}_{/S}\)-groups; or equivalently, the morphism of
  \(\widehat{\mathscr C}\)-groups
  \(\mathbf G\to\underline{\operatorname{Aut}}(\mathbf H)\) maps
  \(\mathbf G\) into
  \[
    \underline{\operatorname{Aut}}_{
      \widehat{\mathscr C}\text{-}\mathbf{Gr}}(\mathbf H).
  \]

  \item[(b)]
  In the situation above, there is a unique
  \(\widehat{\mathscr C}\)-group structure on
  \(\mathbf H\times\mathbf G\) such that, for every \(S\),
  \((\mathbf H\times\mathbf G)(S)\) is the semidirect product of the
  groups \(\mathbf H(S)\) and \(\mathbf G(S)\) for the given action of
  \(\mathbf G(S)\) on \(\mathbf H(S)\). We denote this
  \(\widehat{\mathscr C}\)-group by \(\mathbf H\mathbin{\cdot}\mathbf G\)
  and call it the \emph{semidirect product of \(\mathbf H\) by
  \(\mathbf G\)}. Thus, by definition,
  \[
    (\mathbf H\mathbin{\cdot}\mathbf G)(S)
      =\mathbf H(S)\mathbin{\cdot}\mathbf G(S).
  \]
\end{enumerate}

Let \(\mathbf G\) be a \(\widehat{\mathscr C}\)-group. For every arrow
\(S'\to S\) in \(\mathscr C\) and every \(g\in\mathbf G(S)\), let
\(\operatorname{Int}(g)\) be the automorphism of \(\mathbf G(S')\)
defined by
\[
  \operatorname{Int}(g)h=ghg^{-1}.
\]
This definition extends to a morphism of
\(\widehat{\mathscr C}\)-groups
\[
  \operatorname{Int}:\mathbf G\longrightarrow
  \underline{\operatorname{Aut}}_{
    \widehat{\mathscr C}\text{-}\mathbf{Gr}}(\mathbf G)
  \subset
  \underline{\operatorname{Aut}}(\mathbf G).
\]

\SGARefLink{sga3:i:definition:2:3:3}{Definition 2.3.3} therefore applies, and for every subobject
\(\mathbf E\) of \(\mathbf G\) one has the sub-
\(\widehat{\mathscr C}\)-groups of \(\mathbf G\)
\[
  \underline{\operatorname{Norm}}_{\mathbf G}(\mathbf E)
  \qquad\text{and}\qquad
  \underline{\operatorname{Centr}}_{\mathbf G}(\mathbf E).
\]

\medskip
\SGARefTarget{sga3:i:definition:2:3:6}
\noindent\textbf{Definition 2.3.6.}
The \emph{center} of \(\mathbf G\), denoted
\(\underline{\operatorname{Centr}}(\mathbf G)\), is the
sub-\(\widehat{\mathscr C}\)-group
\(\underline{\operatorname{Centr}}_{\mathbf G}(\mathbf G)\) of
\(\mathbf G\). We say that \(\mathbf G\) is commutative if
\(\underline{\operatorname{Centr}}(\mathbf G)=\mathbf G\), or,
equivalently, if \(\mathbf G(S)\) is commutative for every \(S\).

We say that the sub-\(\widehat{\mathscr C}\)-group \(\mathbf H\) of
\(\mathbf G\) is invariant in \(\mathbf G\) if
\(\underline{\operatorname{Norm}}_{\mathbf G}(\mathbf H)=\mathbf G\),
or, equivalently, if \(\mathbf H(S)\) is invariant in \(\mathbf G(S)\)
for every \(S\).

\par\medskip\noindent\SGAInlineRefTarget{sga3:i:definition:2-3-6-1}%
\emph{Addition} (after the Gille--Polo re-edition (2011), Exposé~I, 2.3.6.1
and 2.3.6.2, and a note to 2.3.6; our formulation). The \emph{kernel} of a
morphism of \(\widehat{\mathscr C}\)-groups \(f:\mathbf G\to\mathbf G'\) is the
sub-\(\widehat{\mathscr C}\)-group \(\operatorname{Ker}f\) of \(\mathbf G\)
with \((\operatorname{Ker}f)(S)=\operatorname{Ker}(f(S))\); it is invariant in
\(\mathbf G\). It is the fibre product of \(f\) and of the unit section of
\(\mathbf G'\); so, if \(\mathscr C\) has a final object \(S_0\) and fibre
products and \(f\) comes from a morphism of \(\mathscr C\)-groups
\(G\to G'\), it is represented by \(S_0\times_{G'}G\). An action of
\(\mathbf G\) on an object \(\mathbf E\) is \emph{faithful} if the kernel of
\(\rho:\mathbf G\to\underline{\operatorname{Aut}}(\mathbf E)\) is trivial,
that is, if an element \(g\in\mathbf G(S)\) acting trivially on every
\(\mathbf E(S')\), \(S'\to S\), is the unit element. A
sub-\(\widehat{\mathscr C}\)-group \(\mathbf H\) of \(\mathbf G\) is
\emph{central} if
\(\underline{\operatorname{Centr}}_{\mathbf G}(\mathbf H)=\mathbf G\), that
is, if \(\mathbf H(S)\) is central in \(\mathbf G(S)\) for every \(S\).
\par\medskip

Many definitions and propositions from elementary group theory carry
over readily. We record only the following one, which will be useful to
us.

\medskip
\noindent\SGAInlineRefTarget{sga3:i:proposition:2-3-7}\textbf{Proposition 2.3.7.}
\textit{Let \(f:\mathbf W\to\mathbf G\) be a morphism of
\(\widehat{\mathscr C}\)-groups. Put
\(\mathbf H(S)=\operatorname{Ker}f(S)\). Let
\(u:\mathbf G\to\mathbf W\) be a morphism of
\(\widehat{\mathscr C}\)-groups that is a section of \(f\) (and is
therefore necessarily a monomorphism). Then \(\mathbf W\) is identified
with the semidirect product of \(\mathbf H\) by \(\mathbf G\) for the
action of \(\mathbf G\) on \(\mathbf H\) defined by
\[
  (g,h)\longmapsto\operatorname{Int}(u(g))h
\]
for \(g\in\mathbf G(S)\), \(h\in\mathbf H(S)\), and
\(S\in\operatorname{Ob}(\mathscr C)\).}

All these definitions and propositions carry over, as usual, to
\(\mathscr C\). In particular, the semidirect product of two
\(\mathscr C\)-groups \(H\) and \(G\), with \(G\) acting on \(H\), is
defined whenever the Cartesian product \(H\times G\) exists, and the
analog of \SGARefLink{sga3:i:proposition:2-3-7}{Proposition 2.3.7} takes the following form.

\medskip
\noindent\SGAInlineRefTarget{sga3:i:proposition:2-3-8}\textbf{Proposition 2.3.8.}
\textit{Let
\[
  H\xrightarrow{i}W\xrightarrow{f}G
\]
be a sequence of morphisms of \(\mathscr C\)-groups such that, for every
\(S\in\operatorname{Ob}(\mathscr C)\), the pair \((H(S),i(S))\) is a
kernel of \(f(S):W(S)\to G(S)\). Let \(u:G\to W\) be a morphism of
\(\mathscr C\)-groups that is a section of \(f\). Then \(W\) is
identified with the semidirect product of \(H\) by \(G\) for the action
of \(G\) on \(H\) such that, if \(S\in\operatorname{Ob}(\mathscr C)\),
\(g\in G(S)\), and \(h\in H(S)\), then
\[
  \operatorname{Int}(u(g))i(h)=i({}^{g}h).
\]}
